\documentclass[11pt]{article}
\usepackage[top=.76in,bottom=.76in,left=.85in,right=.85in]{geometry}
\usepackage{amsmath,amssymb,amsthm,booktabs,microtype,xcolor}
\usepackage[colorlinks=true,urlcolor=blue!45!black]{hyperref}
\newtheorem{theorem}{Theorem}
\newcommand{\F}{\mathbb F}
\title{Eleven is the only sharp prime\\for the universal five-step bound}
\author{}
\date{27 November 2026}
\begin{document}
\maketitle
\vspace{-2em}
\begin{abstract}
Among rank-two Cayley graph towers with five steps, the universal
spanning-tree valuation bound is attained only at $p=11$, and
the maximizing residue list there is unique up to the natural
symmetries. The all-primes exclusion uses 384 cyclotomic norms of
degree four. The same certificate classifies maximizers at $p=31$;
a separate search gives the exact maximum at $p=131$.
\end{abstract}
\section{Statement and reduction}
Let $p>10$ be prime. Let $S=(s_0,\ldots,s_4)$ be a list of integer
vectors spanning $\F_p^2$ after reduction, with repetitions allowed.
For the Cayley multigraph on $(\mathbb Z/p^n\mathbb Z)^2$ with one
undirected edge from $v$ to $v+s_i$ for each $(v,i)$, write
$\kappa_n(S)$ for the spanning-tree count. As in the
\href{https://github.com/VladimirReshetnikov/ProveIt/blob/21c0252ffe8611baa0fbf528fbb003d6c42ca9af/Analysis/GeneralizedRiemannHypothesis/Research/grh_transfers_lubin_tate_and_graph_towers/graphs.tex}{ProveIt graph-tower report}, put
\[
 M_j(u)=\sum_i(u\cdot s_i)^{2j},\quad
 r(L)=\min\{j\ge1:M_j(L)\ne0\},\quad
 A=2\sum_{L\in\mathbb P^1(\F_p)}r(L).
\]
The known moment identity is $v_p(\kappa_n(S))=A(p^n-1)/(p-1)-2n$,
so optimizing $A$ is equivalent to optimizing every layer.
\begin{theorem}
The upper bound $A\le4(p+1)+24$ is attained if and only if $p=11$.
At that prime every maximizer is equivalent to
$((1,0),(0,1),(0,2),(1,1),(3,7))$ under a common invertible linear
map on the residues, permutation, and independent signs.
At $p=31$ the maximum is $148$, also in one orbit, represented by
$((1,1),(2,29),(4,8),(8,16),(16,27))$.
For all other primes $p>10$, $A\le4(p+1)+18$, and if
$p\not\equiv1\bmod 10$ this improves to $4(p+1)+16$.
\end{theorem}

Newton's identities give $r(L)\le5$ by spanning. When $M_1$ is
nonzero the quadratic has at most two projective roots, so
$A\le2(p+1)+16$. Otherwise $M_2$ is a nonzero quartic by the
\href{https://github.com/VladimirReshetnikov/ProveIt/blob/21c0252ffe8611baa0fbf528fbb003d6c42ca9af/Analysis/GeneralizedRiemannHypothesis/Research/grh_transfer_extensions_2026-11-26/sections/graph_extremum.tex}{first-moment bound}; zero steps may be discarded for that lemma.
The baseline is then $4(p+1)$ with at most four exceptional
directions, each contributing $2(r-2)$. In particular, zero, one,
or four depth-five directions allow respectively at most $16$, $18$,
or $24$ above baseline; attaining $24$ requires four depth-five directions.

If $r(L)=5$, the squared projections have monic polynomial $T^5-c$
with $c\ne0$, by Newton's identities. They are a coset of $\mu_5$
inside the nonzero squares. This implies $p\equiv1\bmod 10$ and also
rules out any zero residue step in such a configuration.
\section{Two depth-five directions force a small certificate}
Take any two depth-five directions as coordinate covectors. After
independent row scalings, each row consists of the five powers of a
fixed primitive fifth root $z\in\F_p$, up to independent coordinate
signs: squaring is a bijection on $\mu_5$.
Permute and sign the columns to put the first row in the order $z^i$,
$i=0,\ldots,4$. Scale the second row so its zeroth entry is $1$.
It now has entries $\epsilon_i z^{\sigma(i)}$, with
$\epsilon_0=1$ and $\sigma(0)=0$. The cross term in $M_1=0$ gives
\[
 f(z)=0,\qquad f(Z)=\sum_{i=0}^4\epsilon_i Z^{(i+\sigma(i))\bmod 5}.
\]
There are exactly $4!\,2^4=384$ choices. Multiplication by $f$ in
$\mathbb Z[Z]/(\Phi_5)$ has matrix with entry
$f_{(i-j)\bmod 5}-f_{(4-j)\bmod 5}$ at $(i,j)$, $0\le i,j<4$,
where $f_h$ is the coefficient of $Z^h$. If $f(z)=0$ in $\F_p$,
evaluation at $z$ annihilates the image of this matrix modulo $p$,
so its determinant is divisible by $p$. The supplied checker computes
the absolute determinants by the four-by-four Leibniz formula:
\[
 \{0,1,11,16,25,31,81,101,121,131,625\}.
\]
Thus for $p>10$ every case with nonzero determinant reduces to
$p\in\{11,31,101,131\}$.
For determinant zero the coefficient vector is constant because
$\Phi_5$ is irreducible over $\mathbb Q$. Its odd coefficient sum,
with $\epsilon_0=1$, forces all five signs to be positive.
The checker gives the three permutations $\sigma(i)=ki\bmod5$
with $k=1,2,3$. For $k=1$ the rows coincide. For $k=2,3$ a geometric
sum gives $M_2=20X^3Y$ or $20XY^3$, respectively; only the axes
are exceptional, both of depth five. This gives excess $12$.

For the four exceptional primes, evaluation of the 384 cases at
one primitive fifth root $z$ per prime gives the following full
list of surviving patterns. The columns start at $(1,1)$, so
dependent rows are exactly coincident rows and are discarded.
Counts refer to normalized cases, not equivalence classes.
\begin{center}
\begin{tabular}{@{}rrl@{}}
\toprule
$p$ & $z$ & Exceptional depths with multiplicities\\
\midrule
11  & 3  & $(5,5)$: 2;\quad $(5,5,5,5)$: 30\\
31  & 2  & $(5,5)$: 2;\quad $(3,5,5,5)$: 10\\
101 & 36 & $(5,5)$: 2;\quad $(3,3,5,5)$: 5\\
131 & 53 & $(5,5)$: 7\\
\bottomrule
\end{tabular}
\end{center}
The table proves that four depth-five directions occur only at $11$.
At $31$ the best excess with two or more depth-five directions
is $20$, while at most one allows only $18$. At the other primes
the two-direction cases give excess at most $16$. This proves all
the numerical bounds in the theorem.

For the orbit check take each ordered independent column pair and
both relative signs, map that pair to the coordinate basis, then
take the smaller representative of each sign pair and sort the list.
The least resulting list is constant on each orbit and identifies it,
since every move is invertible. The 30 sharp cases at $11$ reduce
to the representative in the theorem, and the 10 maximal cases at
$31$ reduce to the orbit of the stated witness.

\paragraph{Independent benchmark at 131.}
The small certificate alone does not exclude scores above $540$
from lists with fewer than two depth-five directions. The second
checker normalizes two independent steps to the basis and enumerates
the three remaining signed classes with zero quadratic moment.
This uses the injective key $(x^2,xy,y^2)$ and determines the last
class from the other two. It gives $A_{\max}=540$ even allowing
zero residues. The witness
$((1,0),(0,1),(0,42),(1,20),(28,18))$ has distinct nonzero
signed classes, hence attains the maximum in simple graphs, with
$v_{131}(\kappa_n)=540(131^n-1)/130-2n$.

\paragraph{Scope.}
The all-primes theorem settles sharpness of the universal bound
for five steps, not the exact maximum for every prime. It concerns
residue lists, not equivalence of all integral lifts, and uses no RH
or GRH assumption.
\end{document}
